\section{Project goals and the results established so far}\label{sec:goals}

The project studies the reciprocal differential on finite covers of the
canonical theta heat field. The intended long-term application is to
control degeneration as zeros approach one another. The distinction between
an identity and an estimate is central. Expressing simplicity as local
integrability does not prove simplicity. A useful advance would be an
upper estimate obtained from the unchanged canonical source and remaining
controlled in the relevant limiting regime, or a demonstrable finite
acquisition advantage over direct integration using the same information.

This manuscript develops the analytic part of that program. Its setting
is a bounded disk containing two simple real zeros. The resulting cover
is an annulus, with two physical boundary components, one absolute core
cycle, and a relative path connecting the two boundary components. The
relative period must be retained: it can diverge even when the core period
is constant and the exterior differential has uniformly bounded traces.

\subsection{Canonical normalization and the motivating question}
We use
\begin{align}
 H(t,z)&=\int_0^\infty e^{tu^2}\Phi(u)\cos(zu)\dd u,\label{eq:Hdef}\\
 \Phi(u)&=\sum_{n=1}^\infty
 (2\pi^2n^4e^{9u}-3\pi n^2e^{5u})e^{-\pi n^2e^{4u}}.
 \label{eq:Phi}
\end{align}
The classical theta identity fixes the normalization
\begin{equation}\label{eq:xi}
 H(0,z)=\frac18\xi\left(\frac12+\frac{\ii z}{2}\right),\qquad
 \xi(s)=\frac{s(s-1)}2\pi^{-s/2}\Gamma(s/2)\zeta(s).
\end{equation}
Appendix~\ref{sec:normalization} derives this identity and the complete
source reflection $\Phi(u)=\Phi(-u)$ from the Gaussian theta identity.
The convention agrees with \cite[(1)--(4)]{RT}.
Our derivative and energy estimates below start from
\eqref{eq:Hdef}--\eqref{eq:Phi}. In particular,
\begin{equation}\label{eq:heat}
 \partial_t H=-\partial_z^2H.
\end{equation}
The RH motivation concerns the zero set of \eqref{eq:xi}; none of the local
theorems below assumes RH. In a hypothetical ordinary collision, the
zeros split at the square-root scale in positive physical time. The
reciprocal energy then grows logarithmically. Preventing such a collision
would require a source estimate incompatible with that growth, with
constants valid in the limiting family. The present finite estimates do
not have that property.

\subsection{Analytic contributions}
The results are organized as follows.
\begin{enumerate}
\item Section~\ref{sec:geometry} constructs the cover, normalizes the
energy, and proves the local integrability statement at all multiplicities.
It gives explicit coordinates for the full bordered annulus.
\item Section~\ref{sec:boundary} proves the complete cut-boundary identity
and a real-primitive reduction. With consistently oriented periods $q,p$,
\begin{equation}\label{eq:roadmapGreen}
 \E=\im(\bar q p)+R_{\ext},\qquad
 |R_{\ext}|\le 2\pi^2K_\partial^2.
\end{equation}
The bound uses the exterior differential trace; the relative period
contains the missing degeneration.
\item Section~\ref{sec:quadratic} evaluates the entire quadratic cover and
each of its cut and physical-boundary contributions. It proves a scoped
failure of upper estimates based only on fixed exterior norm budgets and
the permitted core period.
\item Section~\ref{sec:remainders} subtracts a quadratic differential with
the same core period. It proves explicit bounds for the full energy,
relative period, and both boundary contributions, with constants independent
of the root separation when the unit and its derivatives are controlled.
\item Section~\ref{sec:heat} derives a complete quantitative chart from
canonical source moments at a hypothetical ordinary event. It proves a
finite logarithmic upper bound and a strict improvement over its coarse
version on a specified smaller window. These are shared direct estimates,
not a new restriction on the source's admissible events.
\end{enumerate}
Section~\ref{sec:enclosures} records analytic error consequences. Numerical
support is confined to Section~\ref{sec:support}, with the detailed
computational exposition left to its cited record. Section~\ref{sec:open}
identifies the remaining endpoint and coverage problems.

\subsection{Comparison principle and scope}
The direct comparator is always the same quantity,
$2\int_\Omega |H|^{-1}\dd A$, supplied with all the same source,
normalization, root, boundary, and period observations. It may use the
same coordinate substitutions and exact identities. Consequently, a
better bound than one coarse direct certificate does not establish a
general advantage over direct integration. Likewise, a native numerical
implementation can be faster without having more information.

The results collected here consolidate \cite{Notes}; their proofs are
included in full. Classical complex analysis, the exact canonical
normalization above, and ordinary calculus supply the background. There
is no compact finite-genus model of the complete entire theta cover in
the argument, and no omission of the spatial boundary of a finite cover.
